\begin{papercard}{braun2021}
\cf{Stages} \stagetag{C} \stagetag{I} \stagetag{R}. A methodological critique that \emph{separates three families of thematic analysis} (coding-reliability, codebook, reflexive) and states which quality criteria belong to which. It lists ten common misuses and ends with a 20-question evaluation tool. The decisive paper for choosing and defending the thesis' analytic family and its reliability stance.

\cf{Core claim} Universal quality standards for TA are difficult because TA exists in versions that differ in paradigm, philosophy and procedure. Reviewers often demand ``coding reliability measures'' as if they were universal \pc{braun2021}{1}{Developing a universal quality standard for thematic analysis (TA) is complicated by the existence of numerous iterations of TA}. The origin was a reviewer's request to ``discuss how they attempted to avoid bias''. Kidder \& Fine's \emph{small q} (qualitative techniques inside a positivist paradigm) vs.\ \emph{Big Q} (qualitative paradigm), and ``confused q'' as an unknowing mix \pc{braun2021}{2}{we find Kidder and Fine’s (1987) distinction between small q and Big Q qualitative research useful}.

\cf{Reflexive TA} Six phases: (1) familiarisation and notes; (2) systematic coding; (3) \emph{generating} initial themes; (4) reviewing and refining; (5) defining and naming; (6) writing up. The phases are not rigid and become recursive \pc{braun2021}{4}{data familiarisation and writing familiarisation notes; 2) systematic data coding}. It can be inductive or deductive (a continuum), semantic or latent, and experiential or critical. Researcher subjectivity is an analytic resource. Themes are an \emph{outcome} of coding, which is not a search for pre-set themes \pc{braun2021}{4}{Within reflexive TA, the coding process is integral to theme development}.

\cf{The three TA families (Problem 1)} \textbf{Coding-reliability TA}: neopositivist, concerned with objective, unbiased coding. Uses a codebook, multiple independent coders and inter-rater reliability, with early or deductive themes. \textbf{Codebook TA} (framework, matrix, template analysis): a structured coding frame and early themes, in a qualitative paradigm with pragmatic compromises; consensus and IRR are usually not quality measures. \textbf{Reflexive TA}: no coding framework, no team required, themes as the final outcome \pc{braun2021}{6}{‘Coding reliability’ TA captures neopositivist approaches that have at their core concerns about ‘objective’ and ‘unbiased’ coding}. Applying reliability criteria to reflexive TA is ``illogical, incoherent and ultimately meaningless'' \pc{braun2021}{7}{Demonstrating coding reliability and the avoidance of ‘bias’ is illogical, incoherent and ultimately meaningless in a qualitative paradigm}. Codebook approaches suit applied, team-based, deadline-driven work with descriptive, actionable output \pc{braun2021}{7}{These approaches often demonstrate a ‘qualitative pragmatism’ and work well for applied research}.

\cf{Problems 2--10 in brief} (2) Citing without reading: authors claim to follow Braun \& Clarke (2006) but describe codebooks, consensus coding, IRR and early themes that the 2006 paper does not advocate \pc{braun2021}{8}{We promote the use of codebooks and coding frames, consensus coding, the measurement of coding reliability}. (3) Unjustified mash-ups (reflexive TA plus codebooks and reliability). (4) Assuming TA is atheoretical: even inductive TA needs stated assumptions. (5) Assuming TA is only realist or experiential. (6) Assuming TA is only descriptive. (7) Confusing codes (single-facet analytic units) and themes (multi-faceted patterns). (8) Confusing themes and topics: shared meaning around a central concept vs.\ shared topic. (9) ``Themes emerged'': themes are generated, not discovered, so the phase was relabelled ``generating initial themes'' \pc{braun2021}{14}{we have, for now, relabelled this phase ‘generating initial themes’}. (10) Uncritical acceptance of procedure: ``a compass and a map'', not a holiday package.

\cf{The evaluation tool (Table 1, 20 questions)} Grouped as: adequate choice and explanation of methods and methodology (Q1--13: is TA justified, is the \emph{type} of TA specified, does it fit the aims, theory and data collection, is it consistently enacted, are problematic assumptions or mash-ups present, are supplementary methods justified, are theoretical underpinnings and researcher positioning stated, are the procedures described as actually done, is there conceptual or procedural confusion); and a well-developed analysis (Q14--20: are themes clear, are they topic summaries, is contextual information presented as a theme, are applied themes actionable, conceptual clashes, weak analysis such as too many themes or code/theme confusion, problematic statements about generalisability) \pc{braun2021}{24}{Do the authors clearly specify and justify which type of TA they are using}, \pc{braun2021}{25}{Is there evidence of conceptual and procedural confusion}.

\cf{Use for this thesis} Mandatory for the methodology chapter's positioning. The planned design (a predefined codebook, two independent coders, reliability statistics, counts across apps) is, in this paper's terms, \emph{coding-reliability or codebook TA/content analysis}, not reflexive TA. The thesis must therefore (1) name the approach and justify it against Q1--6; (2) \emph{not} cite Braun \& Clarke 2006 as authority for codebooks, IRR or pre-set themes; (3) avoid ``themes emerged'' language and keep the code/theme/topic vocabulary consistent; (4) state positionality and assumptions even for a deductive design; (5) report procedures as actually carried out (Q11). Use the 20 questions as a pre-submission self-audit. Combine with Braun \& Clarke 2006 only for the generic phase vocabulary, and state what is not taken from it.
\end{papercard}
